% 2.5 — brief en documento/investigacion/2-5-gestion-proyectos-epc.md
\subsection{GESTIÓN DE PROYECTOS EPC}
\subsubsection{DEFINICIÓN Y CARACTERÍSTICAS}

Un proyecto EPC (\textit{engineering, procurement and construction}) es aquel en el que un mismo contratista asume, bajo un único contrato, la ingeniería de diseño, la procura de equipos y materiales y la construcción de la instalación. La forma contractual de referencia es el Libro Plata de \textcite{FIDIC2017}, las condiciones de contrato para proyectos EPC «llave en mano», en las que el contratista responde de manera integral por el diseño y la ejecución y entrega al cliente una instalación lista para operar. \textcite{Galloway2009} describe la posición del contratista desde el punto de vista del propietario: una ventanilla única y a la vez la última parada para todos los costos del proyecto, desde su inicio hasta el cierre.

El rasgo que distingue esta modalidad es la distribución del riesgo. La propia FIDIC declara que el Libro Plata no es adecuado cuando el comitente pretende supervisar de cerca la obra ni cuando existe obra subterránea sustancial que los oferentes no pueden inspeccionar \parencite{FIDIC2017}; en otras palabras, es la forma contractual en la que más riesgo se traslada al contratista, y a cambio el cliente obtiene un precio y un plazo comprometidos con menor intervención de su parte. \textcite{Galloway2009} matiza el atractivo de ese arreglo: aunque el esquema puede ahorrar al propietario sumas importantes en el precio inicial y en órdenes de cambio, no es la solución definitiva que se percibe, porque el tratamiento del cambio sigue siendo el problema peor definido en ese entorno. Para el contratista, cada cambio de alcance que no logra trasladar al precio es una desviación que absorbe.

Las tres fases que dan nombre a la modalidad forman una cadena. \textcite{Yeo2002} proponen para los proyectos EPC un marco de procura que integra la gestión de la cadena de suministro con la cadena crítica derivada de la teoría de restricciones, lo que pone de relieve el lugar de la procura, con sus plazos de entrega y sus dependencias con proveedores y subcontratistas, entre la ingeniería y la construcción. En los sectores donde esta modalidad predomina los resultados son con frecuencia adversos: \textcite{Merrow2011}, a partir de plataformas costa afuera, plantas químicas y proyectos de petróleo, gas y minerales, sostiene que más de la mitad de los grandes proyectos de ingeniería y construcción tienen resultados pobres, con sobrecostos y retrasos, y lo atribuye a una gestión inadecuada, a la disfunción de los equipos, a una débil rendición de cuentas y a una inversión técnica insuficiente en las etapas tempranas.

ISI Mustang ejecuta proyectos de este tipo en gas, petróleo, minería y energía, y ocupa en ellos la posición del contratista que describe la literatura: asume la ingeniería, la procura y la construcción, y con ellas el riesgo de que el costo o el plazo se aparten de lo comprometido. En el ERP nuevo cada proyecto se modela como un centro de costo con presupuesto, moneda, responsable y vigencia propios, sobre el que se imputan horas, compras y viáticos; el portal anterior organiza también su histórico por centros de costo. Esa correspondencia entre la unidad contractual y la unidad de registro es lo que hace posible observar, proyecto por proyecto, el fenómeno que las dos subsecciones siguientes describen: cómo se mide el avance y cómo se manifiestan las desviaciones.

\subsubsection{CERTIFICACIONES Y CONTROL DE AVANCE}

El marco de referencia para el control de avance es la gestión del valor ganado (\textit{earned value management}, EVM). \textcite{Anbari2003} la define como un método que integra el alcance, el costo y el tiempo, y que requiere el monitoreo periódico de los gastos reales y de los logros físicos del alcance; a partir de ellos calcula variaciones de costo y de cronograma junto con índices de desempeño, pronostica el costo a la terminación y señala la posible necesidad de acciones correctivas. El \textcite{PMI2019EVM} lo formula en los mismos términos en su norma vigente: una metodología para integrar alcance, cronograma y recursos, medir de forma objetiva el desempeño y el avance, y pronosticar el resultado del proyecto. \textcite{Fleming2010} describen el ciclo de aplicación: definir el alcance mediante una estructura de desglose del trabajo, fijar una línea base de medición del desempeño, medir periódicamente el valor ganado, calcular variaciones e índices y pronosticar el costo final.

El método se apoya en tres magnitudes acumuladas a una fecha de corte: el valor planificado (PV), presupuesto del trabajo que debía haberse hecho; el valor ganado (EV), presupuesto del trabajo efectivamente realizado; y el costo real (AC), lo gastado en ese trabajo. De su comparación derivan las variaciones e índices del cuadro~\ref{tab:evm-indicadores}, que la séptima edición de la guía PMBOK incluye entre las métricas del dominio de desempeño de medición \parencite{PMI2021}. La dificultad práctica está en el valor ganado, porque exige medir el avance físico; la norma admite para ello varios métodos, entre ellos la fórmula fija, los hitos ponderados, el porcentaje completado y la medición física \parencite{PMI2019EVM}.

\begin{table}[H]
  \centering
  \caption{Variaciones, índices y pronóstico de la gestión del valor ganado}\label{tab:evm-indicadores}
  \begin{tblr}{colspec={X[0.96,l] X[0.7,l] X[1.34,l]}}
    \textbf{Indicador} & \textbf{Cálculo} & \textbf{Lectura} \\
    Variación de costo (CV) & $\mathrm{EV}-\mathrm{AC}$ & Negativa: el trabajo realizado costó más de lo presupuestado \\
    Variación de cronograma (SV) & $\mathrm{EV}-\mathrm{PV}$ & Negativa: se hizo menos trabajo del planificado a la fecha \\
    Índice de desempeño de costo (CPI) & $\mathrm{EV}/\mathrm{AC}$ & Menor que 1: sobrecosto en lo ejecutado \\
    Índice de desempeño de cronograma (SPI) & $\mathrm{EV}/\mathrm{PV}$ & Menor que 1: atraso respecto del plan \\
    Estimación a la terminación (EAC) & $\mathrm{BAC}/\mathrm{CPI}$ (forma básica) & Costo final proyectado si se mantiene el desempeño actual; BAC es el presupuesto total \\
  \end{tblr}
  \fuente{Elaboración propia, 2026, en base a \textcite{Anbari2003}, \textcite{Fleming2010} y \textcite{PMI2019EVM}.}
\end{table}

En un contrato EPC el avance es la base del pago. La cláusula de precio y pago del Libro Plata establece que el contratista presenta estados de cuenta periódicos con el monto que considera devengado y que el propio comitente los evalúa y paga, sin el ingeniero independiente que en otras formas de la suite emite el certificado de pago provisional \parencite{FIDIC2017}. En la interpretación de este trabajo, la certificación de avance es el acto contractual por el cual el cliente reconoce como ejecutada una porción del alcance y habilita su facturación. En el ERP nuevo se modela como certificación, el ingreso ejecutado de un centro de costo, que el gerente del centro solicita y el área de compras registra y cobra, con estados de borrador, enviada y registrada. Con la misma lectura, esa entidad cumple la función del ciclo de estado de cuenta, certificación y pago de FIDIC: cada certificación es a la vez una medición de lo ejecutado y un ingreso reconocido.

Los índices del valor ganado se usan además como predictores, según se vio en la sección~2.4.2 con \textcite{Vandevoorde2006} y \textcite{Lipke2009}. Ese uso es el antecedente directo del enfoque de este proyecto, con una precisión de alcance: el ERP nuevo no gestiona una estructura de desglose del trabajo, por lo que no aplica EVM en sentido estricto; lo que ofrece, por centro de costo, es un presupuesto con línea base y seguimiento, los costos imputados y la serie de certificaciones, con los que se prevé aproximar las variaciones e índices del cuadro anterior como variables del modelo, y, como los modelos se entrenan con el histórico del portal anterior, esas variables deben poder reconstruirse también en su esquema, lo que se verifica en el análisis de datos del Capítulo III\@.

\subsubsection{DESVIACIONES PRESUPUESTARIAS Y DE CRONOGRAMA}

Los sobrecostos y los retrasos en grandes proyectos son sistemáticos y no excepcionales. \textcite{Flyvbjerg2014} lo resume en la «ley de hierro» de los megaproyectos: por encima del presupuesto, por encima del plazo, una y otra vez; estima el gasto mundial en ellos entre seis y nueve billones de dólares anuales, cerca del ocho por ciento del producto bruto mundial, y explica la persistencia del patrón por un modelo de romper y arreglar, en el que los proyectos se corrigen sobre la marcha. En la literatura sobre proyectos EPC la evidencia coincide: \textcite{EPCleadingIndicators2018} señalan que más de la mitad de los proyectos enfrenta retrasos significativos y escaladas de costo importantes, y en hidrocarburos \textcite{Olaniran2015} recogen un reporte según el cual el 64\,\% de los megaproyectos en ejecución en el mundo presenta sobrecostos, y advierten que la investigación sobre por qué y cómo ocurren es limitada.

Las causas identificadas se repiten entre estudios. \textcite{Assaf2006}, en una encuesta en Arabia Saudita, identifican 73 causas de retraso, encuentran que el 70\,\% de los proyectos considerados sufrió sobreplazo y que la causa más frecuente para las tres partes es la orden de cambio. \textcite{Olawale2010} distinguen las causas de las desviaciones de los factores que impiden controlarlas: en 250 organizaciones del Reino Unido, los cinco inhibidores principales del control de costo y plazo son los cambios de diseño, los riesgos e incertidumbres, la evaluación inexacta de la duración, la complejidad y el incumplimiento de los subcontratistas; sus medidas de mitigación se clasifican en preventivas, predictivas, correctivas y organizacionales. Para EPC, \textcite{EPCleadingIndicators2018} revisan más de doscientos artículos en busca de indicadores adelantados de costo y plazo por fase de ingeniería, procura y construcción; confirman el cambio de diseño como causa principal, agregan la escasez de recursos y la fluctuación de precios, y observan que todo retraso termina generando sobrecosto.

En el sector de ISI Mustang los estudios son menos numerosos pero apuntan en la misma dirección. \textcite{Olaniran2015} sostienen que los sobrecostos en petróleo y gas surgen de interacciones complejas entre las características del proyecto, las personas, la tecnología, la estructura y la cultura, y proponen la teoría del caos como marco para explicarlas. Un estudio preliminar en fabricación EPC para petróleo y gas en Malasia, basado en entrevistas a diez especialistas, vincula los sobrecostos con la asignación de recursos y ubica en el presupuesto insuficiente el problema principal \parencite{CostOverrunsEPC2022}; su alcance es limitado. El rasgo común de estos trabajos, y de \textcite{Merrow2011}, es que la desviación no obedece a una causa aislada sino a la combinación de muchas, lo que argumenta a favor de un modelo que integre múltiples variables en lugar de una regla de umbral sobre un solo indicador.

Para este proyecto, la literatura reseñada delimita qué se entiende por desviación. La desviación presupuestaria es la diferencia entre el costo real de un proyecto y su presupuesto, que en la nomenclatura del valor ganado corresponde a la variación y al índice de costo; la desviación de cronograma es la diferencia entre la fecha de cierre efectiva y la comprometida; y la desviación en certificaciones es la diferencia entre lo certificado a una fecha y lo planificado certificar, la forma en que el atraso o el sobrecosto se hacen visibles en el ingreso; como problema de clasificación, la etiqueta se obtiene aplicando un umbral a esa diferencia, y la disponibilidad de lo planificado certificar en el histórico se confirma en el Capítulo III\@. Son las tres respuestas que el módulo predictivo busca anticipar. Los indicadores adelantados de \textcite{EPCleadingIndicators2018}, las órdenes de cambio de \textcite{Assaf2006} y los inhibidores de \textcite{Olawale2010} orientan la búsqueda de variables predictoras en el histórico; la categoría de medidas predictivas de estos últimos autores es el lugar que ocupa un módulo de predicción integrado al ERP\@.